\section{Taller resuelto: Respuesta lineal por dos integradores conceptuales}

\noindent\textbf{Nivel:} Intermedio.

\apartado{1. Enunciado}{Para la matriz sintética A=[[-0.7,1],[-3.8,-1.1]] y perturbación inicial (alpha,q)=(0.04,0), calcule autovalores y estado a 12 s. Compare exponencial matricial y una integración de ODE.}

\apartado{2. Datos}{alpha en rad, q en rad/s; matriz no identificada de una aeronave.}

\apartado{3. Incógnitas}{Las magnitudes y comprobaciones solicitadas en el caso.}

\apartado{4. Sistema de referencia}{Estados longitudinales perturbados; tiempo en segundos.}

\apartado{5. Hipótesis físicas}{Sistema lineal autónomo, sin mando y sin saturación.}

\apartado{6. Modelo matemático}{Respuesta lineal por dos integradores conceptuales}

\apartado{7. Ecuaciones simbólicas}{Se parte de los balances y definiciones siguientes.}

\begin{equation}\dot{\vect x}=A\vect x\end{equation}

\begin{equation}\vect x(t)=e^{At}\vect x(0)\end{equation}

\apartado{8. Desarrollo algebraico}{Las operaciones se realizan en este orden, conservando las condiciones del modelo.}

\begin{enumerate}

\item Se comprueba que ambas partes reales son negativas.

\item Se evalúa la exponencial de matriz y se integra la ODE con las mismas condiciones.

\item Se compara el estado completo y no sólo un cruce de alpha por cero.

\item La diferencia numérica mide consistencia de estos algoritmos sobre el modelo lineal.

\end{enumerate}

\apartado{9. Sustitución numérica}{Se conservan los resultados intermedios sin redondear y se muestran todas las operaciones escalares disponibles.}

\textbf{Exponencial desarrollada para esta matriz.} El polinomio característico es $\lambda^2+1.8\lambda+4.57=0$, de donde $\lambda=-0.9\pm i\sqrt{3.76}$. Escribir $A=-0.9I+B$ da $B=\bigl(\begin{smallmatrix}0.2&1\\-3.8&-0.2\end{smallmatrix}\bigr)$ y verificar por multiplicación $B^2=-3.76I$. Agrupar términos pares e impares de la serie exponencial produce:\[\fitmath{e^{At}=e^{-0.9t}\left[I\cos(\sqrt{3.76}t)+\frac{B}{\sqrt{3.76}}\sin(\sqrt{3.76}t)\right]}\]\[\fitmath{\alpha(t)=0.04e^{-0.9t}\left[\cos(\sqrt{3.76}t)+\frac{0.2}{\sqrt{3.76}}\sin(\sqrt{3.76}t)\right]}\]\[\fitmath{q(t)=-\frac{0.152}{\sqrt{3.76}}e^{-0.9t}\sin(\sqrt{3.76}t)}\]Sustituir $t=12$ s da $\alpha=-3.163128\times10^{-7}\,\mathrm{rad}$ y $q=1.530889\times10^{-6}\,\mathrm{s^{-1}}$. Para contrastar con integración, se usa el mismo estado inicial y se controla cada componente por separado; una diferencia entre algoritmos no es validación experimental.





\apartado{10. Resultados con unidades}{Los resultados pertenecen al modelo y al intervalo de datos declarados.}

\begin{center}\begin{tabular}{lr}\toprule Magnitud&Resultado\\\midrule

Parte real modal & $-0.9\,\mathrm{s^{-1}}$\\

Parte imaginaria modal & $1.939072\,\mathrm{s^{-1}}$\\

Alpha a 12 s & $-3.163128\times10^{-7}\,\mathrm{rad}$\\

q a 12 s & $1.530889\times10^{-6}\,\mathrm{s^{-1}}$\\

Error de alpha entre métodos & $4.039651\times10^{-14}\,\mathrm{rad}$\\

Error de q entre métodos & $4.606294\times10^{-14}\,\mathrm{s^{-1}}$\\

\bottomrule\end{tabular}\end{center}

\apartado{11. Interpretación}{El estado final pequeño corresponde a una oscilación amortiguada; el resultado pertenece a la matriz sintética.}

\apartado{12. Comprobación independiente y alcance}{La comparación expm frente a $solve_{ivp}$ se ejecuta en el programa; se conservan tolerancias e intervalo.}

\apartado{13. Limitaciones}{La norma de error agrupa dos estados de unidades distintas sólo como diagnóstico numérico; cada componente se revisa además individualmente.}

\par\smallskip\noindent{\footnotesize\textit{Datos y solución originales. Las sustituciones registradas se reproducen con SymPy y Studium; eso no establece validación experimental ni revisión experta.}}
